\documentclass[a4paper,12pt]{article}
\usepackage[T2A]{fontenc}
\usepackage[utf8]{inputenc}
\usepackage[russian]{babel}
\usepackage{amsmath,amssymb,mathtools}
\usepackage{helvet}
\renewcommand{\familydefault}{\sfdefault}
\usepackage{icomma}
\usepackage[margin=23mm,top=24mm,bottom=23mm]{geometry}
\usepackage{microtype}
\usepackage{tikz}
\usetikzlibrary{calc,arrows.meta,positioning,shapes.geometric}
\usepackage{pgfplots}
\pgfplotsset{compat=1.18}
\usepackage{booktabs,tabularx,longtable}
\usepackage{enumitem}
\usepackage{hyperref}
\usepackage{parskip}
\usepackage{fancyhdr}
\usepackage{setspace}
\usepackage{xcolor}
\definecolor{accent}{HTML}{267F83}
\definecolor{darkaccent}{HTML}{23565A}
\definecolor{textgray}{HTML}{333333}
\definecolor{softgray}{HTML}{6B7475}
\hypersetup{colorlinks=true,linkcolor=darkaccent,urlcolor=darkaccent}
\setstretch{1.17}
\setlength{\parindent}{0pt}
\setlength{\parskip}{0.45em}
\setlist[itemize]{leftmargin=5mm,itemsep=3pt,topsep=3pt}
\setlist[enumerate]{leftmargin=7mm,itemsep=6pt,topsep=4pt}
\setcounter{secnumdepth}{0}
\fancyhf{}
\fancyfoot[L]{\footnotesize\color{softgray}Математика \textbullet\ Графики с модулем}
\fancyfoot[R]{\footnotesize\color{softgray}\thepage}
\renewcommand{\headrulewidth}{0pt}
\pagestyle{fancy}
\newcommand{\accentsection}[1]{\section*{\color{darkaccent}#1}}
\newcommand{\checkitem}[1]{\item[$\square$] #1}
\begin{document}

% Титульный лист
\thispagestyle{empty}
\begin{titlepage}
\centering
\vspace*{35mm}
{\color{accent}\rule{46mm}{1.1pt}}\par
\vspace{21mm}
{\fontsize{23}{28}\selectfont\bfseries ЧЕК-ЛИСТ\par}
\vspace{10mm}
{\fontsize{19}{24}\selectfont\color{darkaccent}\bfseries Графики функций\\ с модулем\par}
\vspace{12mm}
{\large Параболы, гиперболы и пересечения\par}
\vspace{15mm}
{\color{softgray}10 октября 2026 года\par}
\vfill
{\small Лёвин Артём Александрович\\[2mm] эксклюзивно для Марины Селиверстовой\par}
\vspace{15mm}
\begin{tikzpicture}[x=1cm,y=1cm]
  \draw[accent!65,line width=0.9pt] (-6.5,0.4) .. controls (-3.5,-0.4) and (-1.2,-0.25) .. (0,0.0)
      .. controls (1.6,0.28) and (3.5,0.75) .. (6.5,-0.15);
\end{tikzpicture}
\vspace{5mm}
\end{titlepage}

\accentsection{Что необходимо уметь}
\begin{itemize}[label=$\square$]
\item Отличать модуль всей функции от модуля отдельного выражения.
\item Строить исходные параболу и гиперболу по опорным точкам.
\item Отражать относительно оси $Ox$ только участки, где $f(x)<0$.
\item Учитывать знак минус перед всем модулем.
\item Сохранять область определения исходной функции.
\item Определять число пересечений с прямой $y=m$.
\end{itemize}

\accentsection{1. Основное правило}
По определению модуля
\[
|a|=\begin{cases}a,&a\ge 0,\\-a,&a<0.\end{cases}
\qquad
|f(x)|=\begin{cases}f(x),&f(x)\ge0,\\-f(x),&f(x)<0.\end{cases}
\]

Чтобы построить $y=|f(x)|$, сначала постройте $y=f(x)$.
\begin{itemize}
\item Участки \textbf{на оси $Ox$ и над ней} сохраняются.
\item Участки \textbf{под осью $Ox$} отражаются вверх симметрично этой оси.
\item Участки исходного графика под осью удаляются.
\end{itemize}

\textbf{Следствия:} $|f(x)|\ge0$; нули функции сохраняются; область определения остаётся той же, что у $f$.

Если перед модулем стоит минус, $y=-|f(x)|$, полученный график целиком отражают относительно $Ox$. Тогда $y\le0$.

\accentsection{2. Парабола: подробный пример}
Построим график
\[
y=|x^2-x-6|.
\]
Сначала рассмотрим $f(x)=x^2-x-6$:
\[
x_0=-\frac{b}{2a}=\frac12=0{,}5,
\quad f(0{,}5)=-6{,}25,
\quad x^2-x-6=(x+2)(x-3).
\]
Опорные точки исходной параболы:
\[
(0{,}5;-6{,}25),\quad(-2;0),\quad(3;0),\quad(0;-6).
\]
После применения модуля точки с отрицательными ординатами переходят в
\[
(0{,}5;6{,}25),\quad(0;6),
\]
а нули $(-2;0)$ и $(3;0)$ остаются на месте.

\begin{center}
\begin{tikzpicture}
\begin{axis}[
 width=0.94\linewidth,height=72mm,
 axis lines=middle,xmin=-3.4,xmax=4.1,ymin=-7.3,ymax=9,
 axis line style={->,line width=0.65pt},
 xlabel={$x$},ylabel={$y$},
 xtick={-3,-2,-1,0,1,2,3,4},
 ytick={-6,-4,-2,0,2,4,6,8},
 tick label style={font=\scriptsize},label style={font=\small},
 grid=both,major grid style={gray!18},minor grid style={gray!9},
 clip=true,unbounded coords=jump,
 legend style={at={(0.02,0.98)},anchor=north west,draw=none,fill=white,font=\scriptsize}
]
\addplot[domain=-3.2:4,samples=201,gray!75,densely dashed,line width=0.9pt]{x*x-x-6};
\addlegendentry{$y=f(x)$}
\addplot[domain=-3.2:4,samples=201,accent,line width=1.4pt]{abs(x*x-x-6)};
\addlegendentry{$y=|f(x)|$}
\addplot[only marks,mark=*,mark size=1.7pt,color=darkaccent] coordinates {(-2,0) (3,0) (0.5,6.25)};
\end{axis}
\end{tikzpicture}
\end{center}
\textcolor{softgray}{\small Пунктиром показана исходная парабола; сплошной линией --- график с модулем.}

\accentsection{3. Гипербола и внешний минус}
Для $x\ne0$:
\[
\left|\frac1x\right|=\frac1{|x|},\qquad
\left|\frac{-2}{x}\right|=\frac2{|x|}>0,\qquad
-\left|\frac1x\right|=-\frac1{|x|}<0.
\]
\textbf{Пример:} для $y=\left|\dfrac{-2}{x}\right|$ возьмём по три точки с каждой стороны от оси $Oy$.
\begin{center}
\renewcommand{\arraystretch}{1.28}
\begin{tabular}{@{}c*{6}{>{\centering\arraybackslash}p{17mm}}@{}}
\toprule
$x$ & $-4$ & $-2$ & $-1$ & $1$ & $2$ & $4$\\
\midrule
$y$ & $\tfrac12$ & $1$ & $2$ & $2$ & $1$ & $\tfrac12$\\
\bottomrule
\end{tabular}
\end{center}
Обе ветви находятся \textbf{выше} оси $Ox$. Если перед модулем стоит минус, обе ветви окажутся \textbf{ниже} $Ox$. В обеих ситуациях $x=0$ исключён: деление на нуль запрещено.

\accentsection{4. Ошибки, которые важно предупреждать}
\begin{itemize}
\item \textbf{Знак коэффициента $b$:} в формуле $x_0=-b/(2a)$ подставляйте $b$ вместе со знаком. Например, при $b=-3$ получаем $x_0=1{,}5$ для $a=1$.
\item \textbf{Дробь:} $\dfrac32=1{,}5$, а $\dfrac12=0{,}5$.
\item \textbf{Область определения:} модуль не делает допустимым $x=0$ в $1/x$.
\item \textbf{Симметрия:} точка $(x;y)$ переходит в $(x;-y)$ при отражении относительно $Ox$; абсцисса остаётся прежней.
\end{itemize}

\clearpage
\thispagestyle{fancy}
\begin{center}
{\Large\bfseries\color{darkaccent}Алгоритм построения $y=\pm|f(x)|$}
\end{center}
\vspace{8mm}
\begin{center}
\begin{tikzpicture}[
 term/.style={draw=accent,ellipse,align=center,minimum width=38mm,minimum height=10mm,fill=accent!6},
 process/.style={draw=darkaccent,rounded corners=2pt,align=center,text width=56mm,minimum height=14mm,fill=white},
 condition/.style={draw=darkaccent,diamond,aspect=2.3,align=center,text width=30mm,inner sep=1pt,fill=accent!5},
 side/.style={draw=darkaccent,rounded corners=2pt,align=center,text width=41mm,minimum height=17mm,fill=white},
 arr/.style={-{Latex[length=2.5mm]},line width=0.9pt,color=darkaccent},
 every node/.style={font=\small}
]
\node[term] (start) at (0,0) {Начало};
\node[process] (domain) at (0,-1.6) {Найдите область определения $f(x)$};
\node[process] (base) at (0,-3.45) {Постройте исходный график $y=f(x)$};
\node[condition] (q1) at (0,-5.65) {Есть части ниже $Ox$?};
\node[side] (yes1) at (4.3,-8.0) {Отразите нижние части вверх};
\node[side] (no1) at (-4.3,-8.0) {Сохраните график без изменений};
\node[process] (merge) at (0,-10.5) {Получен график $y=|f(x)|$};
\node[condition] (q2) at (0,-12.8) {Перед модулем стоит $-$?};
\node[side] (yes2) at (4.3,-15.05) {Отразите весь график относительно $Ox$};
\node[side] (no2) at (-4.3,-15.05) {Оставьте $y=|f(x)|$};
\node[term] (finish) at (0,-17.5) {Готово};
\draw[arr] (start) -- (domain);
\draw[arr] (domain) -- (base);
\draw[arr] (base) -- (q1);
\draw[arr] (q1.east) -| node[pos=0.22,above]{да} (yes1.north);
\draw[arr] (q1.west) -| node[pos=0.22,above]{нет} (no1.north);
\draw[arr] (yes1.south) |- (merge.east);
\draw[arr] (no1.south) |- (merge.west);
\draw[arr] (merge) -- (q2);
\draw[arr] (q2.east) -| node[pos=0.22,above]{да} (yes2.north);
\draw[arr] (q2.west) -| node[pos=0.22,above]{нет} (no2.north);
\draw[arr] (yes2.south) |- (finish.east);
\draw[arr] (no2.south) |- (finish.west);
\end{tikzpicture}
\end{center}
\clearpage

\accentsection{5. Задача ОГЭ: сколько пересечений возможно?}
Построим $y=|x^2+4x-5|$ и найдём наибольшее число общих точек с прямой, параллельной оси $Ox$.

Для исходной параболы
\[
f(x)=x^2+4x-5=(x+2)^2-9.
\]
Вершина $(-2;-9)$, нули $x=-5$ и $x=1$. После отражения нижнего участка вершина центральной дуги станет точкой $(-2;9)$.

Любая горизонтальная прямая имеет уравнение $y=m$. Перемещаем её мысленно по высоте и считаем \textbf{различные} точки пересечения.
\begin{center}
\renewcommand{\arraystretch}{1.25}
\begin{tabularx}{0.87\linewidth}{@{}>{\centering\arraybackslash}X>{\centering\arraybackslash}X@{}}
\toprule
Положение прямой & Число пересечений\\
\midrule
$m<0$ & $0$\\
$m=0$ & $2$\\
$0<m<9$ & $4$\\
$m=9$ & $3$\\
$m>9$ & $2$\\
\bottomrule
\end{tabularx}
\end{center}
\textbf{Ответ:} наибольшее число общих точек равно $\boxed{4}$.

\textbf{Как проверить рассуждение:} уравнение $|f(x)|=m$ при $m>0$ распадается на два уравнения $f(x)=m$ и $f(x)=-m$. Когда $0<m<9$, оба имеют по два различных корня.

\textbf{Памятка.} Если исходная парабола пересекает $Ox$ в двух точках, а её вершина лежит ниже оси, график с модулем образует центральную дугу; подходящая горизонтальная прямая может пересечь его четыре раза.

\clearpage
\accentsection{Тренировочный блок \textbf{\textemdash} Путь А}
\textcolor{softgray}{10 заданий. От простых преобразований к задачам ОГЭ. Для заданий на график отмечайте опорные точки и характер ветвей.}
\begin{enumerate}
\item При построении $y=|f(x)|$ исходный график проходит через точки $(-2;-3)$, $(0;0)$, $(1;4)$ и $(3;-2)$. Запишите координаты соответствующих точек нового графика.

\item Постройте $y=|x^2-4|$. Укажите нули и точку, в которую перейдёт вершина исходной параболы.

\item Постройте $y=|x^2-2x-3|$. Найдите нули, значение функции при $x=0$ и точку, в которую перейдёт вершина исходной параболы.

\item Постройте $y=|x^2-6x+5|$. Отметьте нули, пересечение с осью $Oy$ и точку, в которую перейдёт вершина исходной параболы.

\item Для $y=\left|\dfrac{-2}{x}\right|$ заполните таблицу при $x=-4,-2,-1,1,2,4$ и постройте график. Укажите ОДЗ.

\item Постройте $y=-\left|\dfrac3x\right|$. Укажите полуплоскость, в которой расположены ветви, ОДЗ и значения функции при $x=-3,-1,1,3$.

\item Сравните $y=|x^2+2x+2|$ и $y=|x^2+2x-2|$. В каком случае исходную параболу приходится отражать? Объясните по положению её вершины и нулей.

\item Постройте $y=|x^2+2x-8|$. Найдите наибольшее число пересечений с горизонтальной прямой $y=m$ и укажите значения $m$, при которых это число достигается.

\item Для графика $y=|x^2-4x+1|$ найдите все значения $m$, при которых прямая $y=m$ имеет ровно три общие точки с графиком.

\item Для графика $y=|x^2+2x+2|$ определите наибольшее возможное число общих точек с прямой, параллельной $Ox$. Объясните, почему четырёх точек быть не может.
\end{enumerate}

\clearpage
\accentsection{Ответы и ориентиры для самопроверки}
\small
\setlength{\LTpre}{4pt}\setlength{\LTpost}{4pt}
\renewcommand{\arraystretch}{1.27}
\begin{longtable}{@{}p{8mm}p{\dimexpr\linewidth-14mm\relax}@{}}
\toprule
\textbf{№} & \textbf{Ключевые результаты}\\
\midrule
\endfirsthead
\toprule
\textbf{№} & \textbf{Ключевые результаты}\\
\midrule
\endhead
1 & $(-2;3)$, $(0;0)$, $(1;4)$, $(3;2)$.\\
2 & Нули $-2$ и $2$; исходная вершина $(0;-4)$ переходит в $(0;4)$. График $\ge0$.\\
3 & Нули $-1$ и $3$; $y(0)=3$; исходная вершина $(1;-4)$ переходит в $(1;4)$.\\
4 & Нули $1$ и $5$; $(0;5)$; исходная вершина $(3;-4)$ переходит в $(3;4)$.\\
5 & $y$: $\frac12,1,2,2,1,\frac12$ соответственно. $x\ne0$; обе ветви выше $Ox$.\\
6 & $x\ne0$; обе ветви ниже $Ox$; $y(-3)=-1$, $y(-1)=-3$, $y(1)=-3$, $y(3)=-1$.\\
7 & Для $x^2+2x+2=(x+1)^2+1$ отражения нет: вершина $(-1;1)$, нулей нет. Для $x^2+2x-2=(x+1)^2-3$ отражается участок между нулями $-1-\sqrt3$ и $-1+\sqrt3$; центральная вершина становится $(-1;3)$.\\
8 & Нули $-4$ и $2$; вершина исходной параболы $(-1;-9)$. Максимум \textbf{4} пересечения при $0<m<9$.\\
9 & Только $m=3$: исходная вершина $(2;-3)$ после отражения даёт вершину центральной дуги $(2;3)$.\\
10 & Максимум \textbf{2} пересечения при $m>1$. Исходная парабола $(x+1)^2+1$ всюду положительна, модуль ничего не меняет.\\
\bottomrule
\end{longtable}
\normalsize
\end{document}
